
Much recent progress in the numerical simulation of mechanical models can be
related to a trend known as the ``geometrization of mechanics''.
Geometry has always played a central role in the development of
mechanics (Euler, Lagrange, Cosserat, etc.),
but with 
the work of authors like Synge, Abraham and Arnold, these geometric 
developments took a new direction as 
% With every confirmation of the predictive power in Einstein's theory of relativity, 
% 
% Replace following with reference to Einstein's relativity and differential
% geometry.
the work of \cite{ricci1901methodes} began to be used to
characterize the geometric structure of non-physical aspects of
problems, like spaces of possible solutions.
This newly found geometric structure in the abstract mathematics would
eventually prove instrumental in the development of computational
methods. 
The work of Simo \citep{simo1986threedimensional,simo1990stress} marked a significant advance in this
direction, as it directly drew from this abstract geometric structure to
inform an enlightened computational procedure.

%
\hypertarget{sec:apparatus}{%
\section{Mechanical Governing Equations}\label{sec:apparatus}}
Developments are centered around static equilibrium problems %that are
governed by nonlinear differential equations of the form:
\begin{equation}
\operatorname{B}_\chi \bm{\sigma} - \bm{f} = \dot{\bm{\mu}}
\label{eq:strong}
\end{equation}
where \(\operatorname{B}_\chi\) is a nonlinear
differential operator and the generalized stress \(\boldsymbol{\sigma}\), momentum rate \(\dot{\bm{\mu}}\), and body force \(\bm{f}\) are defined on a body \(\mathcal{B}\) placed in space by a
map \(\chi\). 
In general \(\operatorname{B}\) depends nonlinearly on \(\chi\),
but acts linearly on \(\boldsymbol{\sigma}\). Problems with the form of
Equation~(\ref{eq:strong}) describe the equilibrium of continuua, rods, and shells, whose
configurations \(\chi\) may comprise a \emph{nonlinear} space
\(\mathscr{C}\).

This chapter presents the mathematical foundation for later developments. In
\cref{box:mechanical-bvp} the mathematical structure which underpins
\cref{eq:strong} is summarized, and in \cref{sec:derivatives} attention
is brought to some important subtleties that arise with derivatives on manifolds. 
This material is adapted from treatments by \cite{marsden1994mathematical,zienkiewicz2014finite,ogden1997nonlinear} and \cite{shilov1996elementary}.

\begin{ambox}[The Mechanical BVP]{box:mechanical-bvp}
{
\begin{description}
\item[Body]
A body \(\mathcal{B}\) is a set of points \(\boldsymbol{\xi}\), and a placement
\(\mathcal{B}_\chi\) of \(\mathcal{B}\) is a Riemannian manifold
\(\mathcal{B}_\chi\) with boundary \(\partial \mathcal{B}\) and tangent
bundle \(T\mathcal{B}_\chi\) embedded in an ambient Euclidean space
\(\mathcal{E}\). The map
\(\chi\colon \mathcal{B}\rightarrow \mathcal{B}_\chi\) associates the
state of a material point \(\boldsymbol{\xi}\) to the placement
\(\mathcal{B}_\chi\). A reference configuration
\(\chi_0\colon \mathcal{B}\rightarrow \mathcal{B}_{\chi_0}\) is often
implied and the pull-back by the composition \(\chi\circ\chi_0^{-1}\) is
denoted by \(\boldsymbol{A}\).
% \tightlist
\item[Metric]
The set \(\mathscr{C}\) of configurations \(\chi\) forms a smooth
 manifold. A configuration \(\chi\) is identified by a couple
\((\boldsymbol{x}, \boldsymbol{\Lambda})\) where \(\boldsymbol{x}\in\mathscr{C}_x\) is a
vectorial configuration variable, and
\(\boldsymbol{\Lambda}\in\mathscr{C}_{\scriptscriptstyle{\Lambda}}\) a
non-vectorial variable. These enter independently in a Riemannian metric
on the tangent space \(T_\chi\mathscr{C}\) to \(\mathscr{C}\) at
\(\chi\), which takes the assumed form: 
\begin{equation}
  \left< \cdot , \, \cdot\right>_{\chi} = \left< \cdot , \, \cdot\right>_{x} + \left< \cdot , \, \cdot\right>_{\scriptscriptstyle{\Lambda}} .
  \label{eq:conf-pair}
\end{equation}
%
\item[Stress]
The spatial stress \(\boldsymbol{\sigma}\) in Equation~(\ref{eq:strong}) is related to a
material stress \(\boldsymbol{s}\) which is defined with respect to inner products 
in \(\mathscr{C}\) at the current \(\chi\) and reference
\(\chi_0\) configurations:
\begin{equation}
  \left<\boldsymbol{c},\, \boldsymbol{\sigma}\right>_{\chi}
  =\left< \boldsymbol{A}\boldsymbol{c},\, \boldsymbol{s}\right>_{\chi_0}
  = \left<\boldsymbol{c},\, \boldsymbol{A}_*\boldsymbol{s}\right>_{\chi}
  =\left< \boldsymbol{A}\boldsymbol{c},\, \boldsymbol{s}\right>_{\chi_0}\qquad\text{ for any }\boldsymbol{c},
  \label{eq:def-Astar}
\end{equation}
and where the transformation \(\boldsymbol{A}_*\) has been defined.
% with \(\boldsymbol{A}_*\) denoting the Piola transformation. 
Similarly, a dual
\(\operatorname{B}^*\) to the equilibrium operator \(\operatorname{B}\) in
Equation~(\ref{eq:strong}) is defined with respect to Equation~(\ref{eq:conf-pair}) for pure displacement boundary conditions:
\begin{equation}
  \left<\boldsymbol{\eta},\, \operatorname{B}_\chi\boldsymbol{\sigma}\right>_\chi = \left<\operatorname{B}_\chi^*\boldsymbol{\eta},\, \boldsymbol{\sigma}\right>_\chi
  \qquad \text{ for any }\boldsymbol{\eta}
  \label{eq:def-adjoint}
\end{equation}
%
\item[Strain]
The material strain \(\boldsymbol{e}\) is related to 
\(\operatorname{B}^*_\chi\) through the differential relation: 
\begin{equation}
  D\boldsymbol{e}\cdot\boldsymbol{u} \equiv \boldsymbol{A}\operatorname{B}^*_\chi \boldsymbol{u}.
  \label{eq:vstrain}
\end{equation}
%
\item[Constitution]
The constitutive response is defined between the material stress
\(\boldsymbol{s}\) and the material strain \(\boldsymbol{e}\) with material \(\boldsymbol{C}_s\) and
spatial \(\boldsymbol{C}_\sigma\) moduli defined by:
\begin{equation}
  D \boldsymbol{s}\cdot \boldsymbol{u} = \boldsymbol{C}_s \, D \boldsymbol{e}\cdot {\boldsymbol{u}}
  \qquad\text{ and }\qquad
  \boldsymbol{C}_\sigma \triangleq \boldsymbol{A}_*\boldsymbol{C}_s\boldsymbol{A}.
  \label{eq:def-moduli}
\end{equation}
\end{description}
}
\end{ambox}

The structure outlined in Box~\ref{box:mechanical-bvp} 
implies a weak statement
of \cref{eq:strong} that is centered around the Galerkin \emph{residual functional}:
%
\begin{equation}
\mathcal{G}\left[\boldsymbol{\eta}\right] \triangleq \left< \boldsymbol{\eta}, \, \mathbb{B}_\chi \boldsymbol{\sigma}\right>_\chi - \left< \boldsymbol{\eta}, \, \boldsymbol{f}\right>_\chi.
\label{eq:galerkin}
\end{equation}
This informs our interpretation of the \emph{principle of virtual work}: forces are by definition resistance to variations of the configuration, and thus the principle of virtual work mandates that these belong to a vector space that is dual to the tangent space of the configuration manifold.



\begin{exmpl}[Elastostatics.]{ex:elastostatics}
Consider the elastostatic problem:
$$
\operatorname{div}\bm{\sigma} + \bm{b} = \bm{0}
$$
and the tensor products:
$$
\begin{aligned}
& (\mathbf{A} \otimes \mathbf{B})[ \mathbf{S}]=(\mathbf{A} \otimes \mathbf{B}): \mathbf{S} = \mathbf{A} (\mathbf{B}\cdot \mathbf{S}) \\
& (\mathbf{A} \oslash \mathbf{B})[ \mathbf{S}]=(\mathbf{A} \oslash \mathbf{B}): \mathbf{S}=\mathbf{A} (\mathbf{S} \circ \mathbf{B}^{\mathrm{t}})=\mathbf{A} \mathbf{S} \mathbf{B}^{\mathrm{t}} \\
%& (\mathbf{A} ? \mathbf{B}): \mathbf{S}=\mathbf{A} \cdot \mathbf{S}^{\mathrm{t}} \cdot \mathbf{B}^{\mathrm{t}} \\
%& (\mathbf{A} \odot \mathbf{B}): \mathbf{S}=\frac{1}{2}(\mathbf{A} \oslash \mathbf{B}+\mathbf{A} \ominus \mathbf{B}): \mathbf{S}=\frac{1}{2}\left(\mathbf{A} \mathbf{S} \mathbf{B}^{\mathrm{t}}+\mathbf{A} \mathbf{S}^{\mathrm{t}}  \mathbf{B}^{\mathrm{t}}\right)
\end{aligned}
$$
\begin{description}
    \item[Body] 
    For a continuum body with reference configuration 
    $\mathcal{B}_0 \subset \mathbb{R}^{n_{\mathrm{dim}}}$ 
    and prescribed deformations on
    the part of the boundary $\Gamma_{\chi} \subset \partial \Omega$, define the configuration space as the set
    $$
    \mathscr{C}=\left\{\chi \in W^{1, p}(\Omega)^{n_{\mathrm{dim}}}
    \quad|\quad
    \det[D \chi]>0 \text{ in }\Omega\text{ and }\left.\chi\right|_{\Gamma_{\chi}}=\bar{\chi}\right\}
    $$
    $$
    \mathcal{V}_{\chi}=\left\{\bm{\eta}: \bm{\chi}(\mathscr{B}) \rightarrow \mathbb{R}^{n_{\mathrm{dim}}}: \bm{\eta}(\chi (\bm{\xi}))
    =\bm{0}\text{ for }\bm{\xi} \in \Gamma_{\chi}\right\}
    $$
    $$
    \mathcal{V}_0 =
    \left\{\bm{\eta}_0 \quad | \quad 
    W^{1, p}(\Omega)^{n_{\mathrm{dim}}}: \bm{\eta}_0(\bm{\xi})=\bm{0} \text{ for } \bm{\xi} \in \Gamma_{\chi}\right\}
    $$
    
    \item[Stress] 
    $\bm{A}_{*} = J^{-1}\bm{F}\oslash\bm{F}$
    $$
    \mathbb{B} = -\operatorname{div}
    \qquad\text{ and }\qquad
    \mathbb{B}^* = \operatorname{grad}
    $$
\end{description}
\end{exmpl}


\iffalse
\hypertarget{sec:statement}{%
\subsection{Problem Statement}\label{sec:statement}}

Equation~(\ref{eq:strong}) is investigated in the context of the following
mathematical structure:

\begin{description}
% \tightlist
\item[A0]
The body \(\mathcal{B}\) is a set of points \(\boldsymbol{\xi}\), and a placement
\(\mathcal{B}_\chi\) of \(\mathcal{B}\) is a Riemannian manifold
\(\mathcal{B}_\chi\) with boundary \(\partial \mathcal{B}\) and tangent
bundle \(T\mathcal{B}_\chi\) embedded in an ambient Euclidean space
\(\mathcal{E}\). The map
\(\chi\colon \mathcal{B}\rightarrow \mathcal{B}_\chi\) associates the
state of a material point \(\boldsymbol{\xi}\) to the placement
\(\mathcal{B}_\chi\). A reference configuration
\(\chi_0\colon \mathcal{B}\rightarrow \mathcal{B}_{\chi_0}\) is often
implied and the pull-back by the composition \(\chi\circ\chi_0^{-1}\) is
denoted by \(\boldsymbol{A}\).
% \tightlist
\item[A1]
The set \(\mathscr{C}\) of configurations \(\chi\) forms a smooth
 manifold. A configuration \(\chi\) is identified by a couple
\((\boldsymbol{x}, \boldsymbol{\Lambda})\) where \(\boldsymbol{x}\in\mathscr{C}_x\) is a
vectorial configuration variable, and
\(\boldsymbol{\Lambda}\in\mathscr{C}_{\scriptscriptstyle{\Lambda}}\) a
non-vectorial variable. These enter independently in a Riemannian metric
on the tangent space \(T_\chi\mathscr{C}\) to \(\mathscr{C}\) at
\(\chi\) which takes the assumed form: 
\begin{equation}
  \left< \cdot , \, \cdot\right>_{\chi} = \left< \cdot , \, \cdot\right>_{x} + \left< \cdot , \, \cdot\right>_{\scriptscriptstyle{\Lambda}} .
  \label{eq:conf-pair}
\end{equation}
%
\item[A2]
The spatial stress \(\boldsymbol{\sigma}\) in Equation~(\ref{eq:strong}) is related to a
material stress \(\boldsymbol{s}\) which is defined with respect to inner products 
in \(\mathscr{C}\) at the current \(\chi\) and reference
\(\chi_0\) configurations:
\begin{equation}
  \left<\boldsymbol{c},\, \boldsymbol{\sigma}\right>_{\chi}
  =\left< \boldsymbol{A}\boldsymbol{c},\, \boldsymbol{s}\right>_{\chi_0}
  = \left<\boldsymbol{c},\, \boldsymbol{A}_*\boldsymbol{s}\right>_{\chi}
  %=\left< \boldsymbol{A}\boldsymbol{c},\, \boldsymbol{s}\right>_{\chi_0}\qquad\text{ for any }\boldsymbol{c},
  \label{eq:def-Astar}
\end{equation}
for any \(\boldsymbol{c}\).
% with \(\boldsymbol{A}_*\) denoting the Piola transformation. 
Similarly, a dual
\(\operatorname{B}^*\) to the equilibrium operator \(\operatorname{B}\) in
Equation~(\ref{eq:strong}) is defined with respect to Equation~(\ref{eq:conf-pair}):
\begin{equation}
  \left<\boldsymbol{\eta},\, \operatorname{B}_\chi\boldsymbol{\sigma}\right>_\chi = \left<\operatorname{B}_\chi^*\boldsymbol{\eta},\, \boldsymbol{\sigma}\right>_\chi.
  \label{eq:def-adjoint}
\end{equation}
% \tightlist
\item[A3]
The material strain measure \(\boldsymbol{e}\) is compatible with
\(\operatorname{B}^*_\chi\) through its directional derivative: 
\begin{equation}
  D\boldsymbol{e}\cdot\boldsymbol{u} \equiv \boldsymbol{A}\operatorname{B}^*_\chi \boldsymbol{u}.
  \label{eq:vstrain}
\end{equation}
%
\item[A4]
The constitutive response is defined between the material stress
\(\boldsymbol{s}\) and the material strain \(\boldsymbol{e}\) with material \(\boldsymbol{C}_s\) and
spatial \(\boldsymbol{C}_\sigma\) moduli defined by:
\begin{equation}
  D \boldsymbol{s}\cdot \boldsymbol{u} = \boldsymbol{C}_s \, D \boldsymbol{e}\cdot {\boldsymbol{u}}
  \qquad\text{ and }\qquad
  \boldsymbol{C}_\sigma \triangleq \boldsymbol{A}_*\boldsymbol{C}_s\boldsymbol{A}.
  \label{eq:def-moduli}
\end{equation}
\end{description}
\fi 
